% ===========================================================================
%  Chapitre 2 — Dérivabilité
%  PE 2026, composante ANALYSE
% ===========================================================================
\bmtchapitre{Dérivabilité}
  {Étudier la dérivabilité d'une fonction en un point et sur un intervalle ;
   dériver une composée et une fonction réciproque ; exploiter la concavité, le
   théorème des accroissements finis et l'inégalité des accroissements finis.}
  {Analyse \textbullet\ environ 10 heures}

En première, la dérivée servait à trouver les variations d'une fonction : on
calculait $f'$, on étudiait son signe, on remplissait un tableau. C'était un
outil de calcul. Cette année, elle devient un outil de démonstration, et le
changement se joue sur un mot : la dérivabilité n'est plus supposée, elle
s'établit.

La différence se voit dès le premier exemple venu. Une fonction peut être
continue en un point sans y être dérivable ; sa courbe présente alors un
angle, et c'est cet angle que le baccalauréat demande d'étudier presque chaque
année. Les deux derniers théorèmes du chapitre, celui des accroissements finis
et son inégalité, permettront de majorer un écart sans jamais calculer la
fonction — ils reviendront, au chapitre des suites, pour prouver une
convergence.

% ---------------------------------------------------------------------------
\begin{revision}
Ce que la première a mis en place, et dont ce chapitre va se servir.

\begin{enumerate}[label=\textbf{\arabic*.}, leftmargin=*, itemsep=0.35em]
  \item Rappeler la définition du nombre dérivé de $f$ en $a$ à l'aide du taux
        d'accroissement.
  \item Donner l'équation de la tangente à la courbe de $f$ au point
        d'abscisse $a$.
  \item Dériver : $x \mapsto x^{4}$, $x \mapsto \sqrt{x}$,
        $x \mapsto \dfrac{1}{x}$, $x \mapsto \cos x$.
  \item Énoncer les formules de dérivation d'un produit et d'un quotient.
  \item Quel lien y a-t-il entre le signe de $f'$ et les variations de $f$ ?
  \item Une fonction continue en $a$ y admet-elle nécessairement une limite
        finie ? Et la réciproque ?
\end{enumerate}
\end{revision}

\begin{automatismes}
Sans calculatrice, sans brouillon.

\begin{enumerate}[label=\textbf{\alph*)}, leftmargin=*, itemsep=0.25em]
  \item Dérivée de $x \mapsto 3x^{2}-5x+1$
  \item Dérivée de $x \mapsto \dfrac{2}{x}$
  \item Dérivée de $x \mapsto (2x+1)^{2}$
  \item Dérivée de $x \mapsto x\sqrt{x}$ sur $\Rps$
  \item Dérivée de $x \mapsto \sin x + \cos x$
  \item Le taux d'accroissement de $x \mapsto x^{2}$ entre $2$ et $2+h$
  \item Tangente à $y = x^{2}$ au point d'abscisse $1$
  \item Dérivée de $x \mapsto \dfrac{x}{x+1}$
\end{enumerate}
\end{automatismes}

\begin{corrige}[automatismes]
a) $6x-5$ \quad b) $-\dfrac{2}{x^{2}}$ \quad c) $4(2x+1)$ \quad
d) $\dfrac{3\sqrt{x}}{2}$ \quad e) $\cos x - \sin x$ \quad
f) $4+h$ \quad g) $y = 2x-1$ \quad h) $\dfrac{1}{(x+1)^{2}}$.
\end{corrige}

% ===========================================================================
\section{Être dérivable en un point}

\begin{definition}[nombre dérivé]
Soit $f$ définie sur un intervalle $I$ contenant $a$. On dit que $f$ est
\textbf{dérivable en $a$} lorsque le taux d'accroissement
\[
  \frac{f(x)-f(a)}{x-a}
\]
admet une limite \emph{finie} quand $x$ tend vers $a$. Cette limite est
appelée le \textbf{nombre dérivé} de $f$ en $a$ et se note $f'(a)$.
\end{definition}

\bmtmarge{En posant $x = a+h$, le taux s'écrit
$\dfrac{f(a+h)-f(a)}{h}$ et l'on fait tendre $h$ vers $0$. Les deux écritures
sont équivalentes ; choisissez celle qui simplifie le calcul.}

\begin{propriete}[dérivabilité et continuité]
Si $f$ est dérivable en $a$, alors $f$ est continue en $a$.
\end{propriete}

\begin{demonstration}
Pour $x \neq a$, écrivons
$f(x) - f(a) = \dfrac{f(x)-f(a)}{x-a} \times (x-a)$.
Quand $x$ tend vers $a$, le premier facteur tend vers $f'(a)$, qui est un
nombre fini, et le second tend vers $0$. Le produit tend donc vers $0$,
c'est-à-dire $\limx{a} f(x) = f(a)$ : $f$ est continue en $a$.
\end{demonstration}

\begin{avertissement}
La réciproque est fausse, et c'est tout l'intérêt de cette propriété. Une
fonction peut être continue en $a$ sans y être dérivable : le taux
d'accroissement existe alors, mais sa limite est infinie, ou bien elle diffère
selon que l'on arrive par la gauche ou par la droite. Retenez le sens de
l'implication : \emph{dérivable} $\implies$ \emph{continue}, jamais l'inverse.
\end{avertissement}

\subsection{Dérivabilité à gauche, à droite, et demi-tangentes}

\begin{definition}[nombres dérivés à gauche et à droite]
On dit que $f$ est dérivable \textbf{à droite} en $a$ lorsque
$\displaystyle\lim_{x \to a^{+}} \frac{f(x)-f(a)}{x-a}$ existe et est finie ;
cette limite est notée $f'_{d}(a)$. On définit de même $f'_{g}(a)$ avec
$x \to a^{-}$.

La fonction $f$ est dérivable en $a$ si et seulement si $f'_{g}(a)$ et
$f'_{d}(a)$ existent, sont finies \emph{et sont égales}.
\end{definition}

\begin{definition}[point angulé]
Lorsque $f$ est continue en $a$ et que $f'_{g}(a)$ et $f'_{d}(a)$ existent
mais sont différents, la courbe admet en son point d'abscisse $a$ deux
\textbf{demi-tangentes} distinctes, de coefficients directeurs respectifs
$f'_{g}(a)$ et $f'_{d}(a)$. On dit que le point est un \textbf{point angulé}.
\end{definition}

\begin{visualisation}[l'angle que fait la courbe]
\begin{center}
\begin{tikzpicture}
\begin{axis}[
  width = 10.6cm, height = 5.6cm,
  axis lines = middle, xlabel = {$x$},
  xmin = -2.6, xmax = 2.8, ymin = -0.5, ymax = 2.6,
  xtick = {0}, ytick = \empty, xticklabels = {$a$},
  axis line style = {bmtGris},
  tick label style = {font=\footnotesize, color=bmtGris}, clip = false,
]
  \addplot[bmtIndigo, line width=1.3pt, domain=-2.4:0, samples=80] {-0.9*x};
  \addplot[bmtIndigo, line width=1.3pt, domain=0:2.6, samples=80] {0.55*x};
  \addplot[bmtLaterite, dashed, line width=0.8pt, domain=-1.4:1.2, samples=2] {-0.9*x};
  \addplot[bmtVetiver, dashed, line width=0.8pt, domain=-1.2:2.2, samples=2] {0.55*x};
  \filldraw[bmtIndigo] (axis cs:0,0) circle (2pt);
  \node[bmtLaterite, font=\footnotesize, anchor=south] at (axis cs:-.9,2.25) {pente $f'_{g}(a)$};
  \node[bmtVetiver, font=\footnotesize, anchor=east] at (axis cs:2.6,1.65) {pente $f'_{d}(a)$};
\end{axis}
\end{tikzpicture}
\end{center}

Approchez du point par la gauche : la courbe arrive avec une pente négative.
Repartez vers la droite : elle repart avec une pente positive. Les deux
demi-tangentes existent, elles sont bien tracées, mais elles ne se prolongent
pas l'une l'autre. Il n'y a donc pas \emph{une} tangente en ce point, et la
fonction n'y est pas dérivable — alors qu'elle y est parfaitement continue,
puisqu'on n'a pas levé le crayon.
\end{visualisation}

\begin{exemple}[une fonction continue et non dérivable]
Soit $f$ définie sur $\R$ par $f(x) = \abs{x}$. Étudions sa dérivabilité en
$0$.

Le taux d'accroissement en $0$ vaut
$\dfrac{f(x)-f(0)}{x-0} = \dfrac{\abs{x}}{x}$.

\smallskip
\emph{À droite.} Pour $x > 0$, $\abs{x} = x$ donc le taux vaut $1$ :
$f'_{d}(0) = 1$.

\smallskip
\emph{À gauche.} Pour $x < 0$, $\abs{x} = -x$ donc le taux vaut $-1$ :
$f'_{g}(0) = -1$.

\smallskip
Les deux nombres existent mais diffèrent : $f$ n'est pas dérivable en $0$.
Elle y est pourtant continue, car $\limx{0} \abs{x} = 0 = f(0)$. La courbe
présente en $O$ un point angulé, avec deux demi-tangentes de pentes $-1$
et $1$.
\end{exemple}

\begin{exemple}[une tangente verticale]
Soit $f(x) = \sqrt{x}$ sur $\left[0 ; +\infty\right[$. En $0$, le taux
d'accroissement vaut
\[
  \frac{\sqrt{x}-0}{x-0} = \frac{\sqrt{x}}{x} = \frac{1}{\sqrt{x}}
  \xrightarrow[x \to 0^{+}]{} +\infty .
\]
La limite est infinie : $f$ n'est pas dérivable en $0$. La courbe y admet une
demi-tangente \emph{verticale}. Ici encore, la fonction est continue en $0$.
\end{exemple}

\begin{entrainement}[dérivabilité en un point]
\exo[\niveau{1}] Soit $f(x) = x^{2}-3x$. Calculer $f'(2)$ en revenant à la
définition, puis vérifier par la formule de dérivation.

\exo[\niveau{2}] Soit $f$ définie par $f(x) = x^{2}$ si $x \leq 1$ et
$f(x) = 2x-1$ si $x > 1$. Étudier la continuité puis la dérivabilité de $f$
en $1$.

\exo[\niveau{2}] Soit $f$ définie par $f(x) = \abs{x-2}$. Montrer que $f$
n'est pas dérivable en $2$ et préciser les demi-tangentes.

\exo[\niveau{3}] Soit $f$ définie sur $\R$ par $f(x) = x\abs{x}$. Montrer que
$f$ est dérivable en $0$ et calculer $f'(0)$. La courbe présente-t-elle un
point angulé ?
\end{entrainement}

\begin{corrige}
\textbf{1.} $\dfrac{f(x)-f(2)}{x-2} = \dfrac{x^{2}-3x+2}{x-2} = \dfrac{(x-2)(x-1)}{x-2} = x-1 \to 1$.
Donc $f'(2) = 1$, ce que confirme $f'(x) = 2x-3$.

\textbf{2.} $\lim_{x \to 1^{-}} x^{2} = 1$, $\lim_{x \to 1^{+}} (2x-1) = 1$ et
$f(1) = 1$ : continue. Pour la dérivabilité,
$f'_{g}(1) = \lim_{x \to 1^{-}} \dfrac{x^{2}-1}{x-1} = \lim (x+1) = 2$ et
$f'_{d}(1) = \lim_{x \to 1^{+}} \dfrac{2x-2}{x-1} = 2$. Les deux valent $2$ :
$f$ est dérivable en $1$, avec $f'(1) = 2$.

\textbf{3.} Le taux vaut $\dfrac{\abs{x-2}}{x-2}$, égal à $1$ pour $x>2$ et à
$-1$ pour $x<2$. Deux demi-tangentes, de pentes $1$ et $-1$ : point angulé.

\textbf{4.} Le taux vaut $\dfrac{x\abs{x}}{x} = \abs{x}$, qui tend vers $0$ des
deux côtés. Donc $f'(0) = 0$ et il n'y a pas de point angulé : la tangente en
$O$ est l'axe des abscisses.
\end{corrige}

% ===========================================================================
\section{Dériver une composée, dériver une réciproque}

\begin{theoreme}[dérivée d'une fonction composée]
Soit $u$ dérivable sur un intervalle $I$ et $v$ dérivable sur un intervalle
$J$ contenant $u(I)$. Alors $v \circ u$ est dérivable sur $I$ et, pour tout
$x \in I$,
\[
  (v \circ u)'(x) = u'(x) \times v'\bigl(u(x)\bigr).
\]
\end{theoreme}

\bmtmarge{La formule se retient dans l'ordre où on l'applique : on dérive
l'extérieur en gardant l'intérieur intact, puis on multiplie par la dérivée de
l'intérieur.}

\begin{propriete}[les cas les plus fréquents]
Soit $u$ dérivable sur $I$. Alors, sous réserve d'existence,
\[
  \left(u^{n}\right)' = n\,u'\,u^{n-1} \ (n \in \Z), \qquad
  \left(\sqrt{u}\right)' = \frac{u'}{2\sqrt{u}}, \qquad
  \left(\frac{1}{u}\right)' = -\frac{u'}{u^{2}},
\]
\[
  \left(\sin u\right)' = u'\cos u, \qquad \left(\cos u\right)' = -u'\sin u .
\]
\end{propriete}

\begin{exemple}[une composée à deux étages]
Dérivons $f(x) = \sqrt{3x^{2}+1}$ sur $\R$.

On pose $u(x) = 3x^{2}+1$, qui est dérivable sur $\R$ et strictement positive,
et $v(X) = \sqrt{X}$, dérivable sur $\Rps$. Alors $f = v \circ u$ et
\[
  f'(x) = \frac{u'(x)}{2\sqrt{u(x)}} = \frac{6x}{2\sqrt{3x^{2}+1}}
        = \frac{3x}{\sqrt{3x^{2}+1}} .
\]
Le signe de $f'$ est celui de $x$ : $f$ décroît sur $\left]-\infty ; 0\right]$
et croît sur $\left[0 ; +\infty\right[$.
\end{exemple}

\begin{theoreme}[dérivée de la fonction réciproque]
Soit $f$ continue et strictement monotone sur un intervalle $I$ ; elle réalise
donc une bijection de $I$ sur $J = f(I)$, de réciproque $f^{-1}$. Si $f$ est
dérivable en $a \in I$ et si $f'(a) \neq 0$, alors $f^{-1}$ est dérivable en
$b = f(a)$ et
\[
  \left(f^{-1}\right)'(b) = \frac{1}{f'(a)} = \frac{1}{f'\bigl(f^{-1}(b)\bigr)} .
\]
\end{theoreme}

\begin{demonstration}[l'idée du calcul]
Posons $y = f(x)$, de sorte que $x = f^{-1}(y)$ et $a = f^{-1}(b)$. Le taux
d'accroissement de $f^{-1}$ en $b$ s'écrit
\[
  \frac{f^{-1}(y)-f^{-1}(b)}{y-b} = \frac{x-a}{f(x)-f(a)}
  = \frac{1}{\ \dfrac{f(x)-f(a)}{x-a}\ }.
\]
Lorsque $y$ tend vers $b$, la continuité de $f^{-1}$ assure que $x$ tend vers
$a$ ; le dénominateur tend donc vers $f'(a)$, non nul. Le quotient tend par
conséquent vers $\dfrac{1}{f'(a)}$.
\end{demonstration}

\begin{remarque}
Géométriquement, les courbes de $f$ et de $f^{-1}$ sont symétriques par
rapport à la droite d'équation $y = x$. Une symétrie échange les rôles des
axes : une tangente de pente $m$ devient une tangente de pente $\frac{1}{m}$,
et une tangente horizontale devient verticale — d'où la condition
$f'(a) \neq 0$.
\end{remarque}

\begin{entrainement}[composées et réciproques]
\exo[\niveau{1}] Dériver $x \mapsto (3x-2)^{5}$, $x \mapsto \sqrt{x^{2}+4}$ et
$x \mapsto \dfrac{1}{2x+3}$.

\exo[\niveau{2}] Dériver $x \mapsto \cos(3x^{2})$ puis $x \mapsto \sin^{3} x$.

\exo[\niveau{2}] Soit $f(x) = x^{3}+x$ sur $\R$. Montrer que $f$ réalise une
bijection de $\R$ sur $\R$, puis calculer $\left(f^{-1}\right)'(2)$.

\exo[\niveau{3}] Soit $f(x) = \dfrac{2x-1}{x+1}$ sur $\left]-1 ; +\infty\right[$.
Déterminer $f^{-1}$ explicitement, puis vérifier la formule du théorème en
un point de votre choix.
\end{entrainement}

\begin{corrige}
\textbf{1.} $15(3x-2)^{4}$ ; \ $\dfrac{x}{\sqrt{x^{2}+4}}$ ; \
$-\dfrac{2}{(2x+3)^{2}}$.

\textbf{2.} $-6x\sin(3x^{2})$ ; \ $3\sin^{2}x\cos x$.

\textbf{3.} $f'(x) = 3x^{2}+1 > 0$ : $f$ est continue et strictement
croissante sur $\R$, avec $\lim_{-\infty} f = -\infty$ et
$\lim_{+\infty} f = +\infty$ ; c'est donc une bijection de $\R$ sur $\R$.
Comme $f(1) = 2$, on a $f^{-1}(2) = 1$ et
$\left(f^{-1}\right)'(2) = \dfrac{1}{f'(1)} = \dfrac{1}{4}$.

\textbf{4.} $y = \dfrac{2x-1}{x+1} \daccord y(x+1) = 2x-1 \daccord x(y-2) = -1-y$,
d'où $f^{-1}(y) = \dfrac{1+y}{2-y}$. En $x = 1$ : $f(1) = \tfrac{1}{2}$,
$f'(x) = \dfrac{3}{(x+1)^{2}}$ donc $f'(1) = \tfrac{3}{4}$ ; par ailleurs
$\left(f^{-1}\right)'(y) = \dfrac{3}{(2-y)^{2}}$ vaut $\tfrac{4}{3}$ en
$y = \tfrac{1}{2}$. Les deux résultats sont bien inverses l'un de l'autre.
\end{corrige}

% ===========================================================================
\section{Dérivées successives, concavité et convexité}

\begin{definition}[dérivées successives]
Si $f'$ est elle-même dérivable sur $I$, sa dérivée est appelée
\textbf{dérivée seconde} de $f$ et notée $f''$. On définit de proche en proche
$f^{(n)}$, dérivée $n$-ième de $f$, par $f^{(0)} = f$ et
$f^{(n+1)} = \left(f^{(n)}\right)'$.
\end{definition}

\begin{definition}[convexité, concavité]
Soit $f$ deux fois dérivable sur $I$. On dit que $f$ est \textbf{convexe} sur
$I$ lorsque $f''(x) \geq 0$ pour tout $x \in I$, et \textbf{concave} lorsque
$f''(x) \leq 0$ pour tout $x \in I$.

Un point de la courbe où la concavité change de sens — donc où $f''$ s'annule
en changeant de signe — est appelé \textbf{point d'inflexion}.
\end{definition}

\begin{propriete}[lecture géométrique]
Sur un intervalle où $f$ est convexe, la courbe est située \emph{au-dessus} de
chacune de ses tangentes. Sur un intervalle où $f$ est concave, elle est
située \emph{au-dessous}. En un point d'inflexion, la tangente traverse la
courbe.
\end{propriete}

\begin{visualisation}[la tangente au-dessus, puis au-dessous]
\begin{center}
\begin{tikzpicture}
\begin{axis}[
  width = 11cm, height = 5.8cm,
  axis lines = middle, xlabel = {$x$},
  xmin = -2.4, xmax = 2.4, ymin = -2.6, ymax = 2.6,
  xtick = {0}, ytick = \empty, xticklabels = {},
  axis line style = {bmtGris},
  tick label style = {font=\footnotesize, color=bmtGris}, clip = false,
]
  \addplot[bmtIndigo, line width=1.3pt, domain=-1.65:1.65, samples=120] {0.55*x^3};
  \addplot[bmtVetiver, dashed, line width=0.85pt, domain=-1.85:-0.3, samples=2]
    {0.55*3*1.44*(x+1.2) - 0.55*1.728};
  \addplot[bmtLaterite, dashed, line width=0.85pt, domain=0.3:1.85, samples=2]
    {0.55*3*1.44*(x-1.2) + 0.55*1.728};
  \filldraw[bmtIndigo] (axis cs:0,0) circle (1.9pt);
  \node[bmtEncre, font=\footnotesize, anchor=south east] at (axis cs:-.1,0.5)
    {point d'inflexion};
  \node[bmtVetiver, font=\footnotesize, anchor=west] at (axis cs:-1.1,-2.1) {concave};
  \node[bmtLaterite, font=\footnotesize, anchor=east] at (axis cs:1.0,2.0) {convexe};
\end{axis}
\end{tikzpicture}
\end{center}

Suivez la courbe de gauche à droite. À gauche de l'origine elle se creuse vers
le bas et passe sous sa tangente : elle est concave. À droite elle se creuse
vers le haut et passe au-dessus : elle est convexe. À l'origine, exactement,
la tangente traverse la courbe — la concavité s'inverse. C'est ce point, et
lui seul, que l'on appelle point d'inflexion.
\end{visualisation}

\begin{exemple}[chercher les points d'inflexion]
Soit $f(x) = x^{4} - 6x^{2} + 5$ sur $\R$.

On calcule $f'(x) = 4x^{3} - 12x$ puis $f''(x) = 12x^{2} - 12 = 12(x-1)(x+1)$.

Le trinôme $f''$ s'annule en $-1$ et en $1$, et son signe est celui d'un
polynôme du second degré de coefficient dominant positif : positif à
l'extérieur des racines, négatif entre elles.

\medskip
\begin{center}
\begin{tabular}{c|ccccc}
$x$        & $-\infty$ & & $-1$ & & $1$ \hspace{0.9cm} $+\infty$ \\
\hline
$f''(x)$   & & $+$ & $0$ & $-$ & $0$ \hspace{1.1cm} $+$ \\
\hline
concavité  & \multicolumn{2}{c}{convexe} & & concave & \hspace{0.4cm} convexe \\
\end{tabular}
\end{center}
\medskip

La dérivée seconde s'annule en changeant de signe en $-1$ et en $1$ : les
points d'abscisses $-1$ et $1$ sont deux points d'inflexion de la courbe.
\end{exemple}

\begin{erreurclassique}
L'annulation de $f''$ ne suffit pas à conclure à un point d'inflexion : encore
faut-il que $f''$ \emph{change de signe}. Pour $f(x) = x^{4}$, on a
$f''(x) = 12x^{2}$, qui s'annule en $0$ sans jamais devenir négatif : la
courbe reste convexe et $O$ n'est pas un point d'inflexion.
\end{erreurclassique}

\begin{entrainement}[concavité]
\exo[\niveau{1}] Calculer $f''$ pour $f(x) = x^{3}-3x^{2}+1$, puis déterminer
les intervalles de convexité et de concavité.

\exo[\niveau{2}] Déterminer les points d'inflexion de la courbe de
$f(x) = \dfrac{1}{x^{2}+1}$.

\exo[\niveau{2}] Montrer que la courbe de $f(x) = \sqrt{x}$ est concave sur
$\Rps$, et en déduire que pour tout $x > 0$, $\sqrt{x} \leq \dfrac{x+1}{2}$.

\exo[\niveau{3}] Soit $f(x) = x^{3}+ax^{2}+bx+c$. Montrer que la courbe de $f$
admet toujours exactement un point d'inflexion, et donner son abscisse en
fonction de $a$.
\end{entrainement}

\begin{corrige}
\textbf{1.} $f''(x) = 6x-6$ : concave sur $\left]-\infty ; 1\right]$, convexe
sur $\left[1 ; +\infty\right[$, point d'inflexion en $x=1$.

\textbf{2.} $f'(x) = \dfrac{-2x}{(x^{2}+1)^{2}}$ puis
$f''(x) = \dfrac{6x^{2}-2}{(x^{2}+1)^{3}}$, qui s'annule en changeant de signe
pour $x = \pm\dfrac{1}{\sqrt{3}}$ : deux points d'inflexion.

\textbf{3.} $f''(x) = -\dfrac{1}{4x\sqrt{x}} < 0$ : $f$ est concave, donc sa
courbe est sous chacune de ses tangentes. La tangente en $x=1$ a pour équation
$y = \tfrac{1}{2}(x-1)+1 = \dfrac{x+1}{2}$, d'où l'inégalité annoncée.

\textbf{4.} $f''(x) = 6x+2a$, qui s'annule en changeant de signe uniquement
pour $x = -\dfrac{a}{3}$ : un point d'inflexion, et un seul.
\end{corrige}

% ===========================================================================
\section{Les accroissements finis}

Les deux résultats qui suivent relient la variation \emph{globale} d'une
fonction, entre deux points, au comportement \emph{local} de sa dérivée. C'est
ce passage du global au local qui en fait des outils de démonstration.

\begin{theoreme}[théorème des accroissements finis]
Soit $f$ continue sur $\intervalle{a}{b}$ et dérivable sur
$\intervalleo{a}{b}$, avec $a < b$. Alors il existe au moins un réel
$c \in \intervalleo{a}{b}$ tel que
\[
  f(b) - f(a) = (b-a)\, f'(c) .
\]
\end{theoreme}

\begin{visualisation}[une tangente parallèle à la corde]
\begin{center}
\begin{tikzpicture}
\begin{axis}[
  width = 11cm, height = 5.8cm,
  axis lines = middle, xlabel = {$x$},
  xmin = -0.5, xmax = 5.4, ymin = -0.6, ymax = 3.9,
  xtick = {0.6, 1.7315, 4.6}, xticklabels = {$a$, $c$, $b$}, ytick = \empty,
  axis line style = {bmtGris},
  tick label style = {font=\footnotesize, color=bmtGris}, clip = false,
]
  % f(x) = 0,42u^2 - 0,06u^3 + 0,5 avec u = x - 0,6 : f(a) = 0,5 et f(b) = 3,38.
  % La corde a donc pour pente 0,72 ; l'equation f'(c) = 0,72 donne c = 1,7315,
  % ou la tangente est bien parallele a la corde.
  \addplot[bmtIndigo, line width=1.3pt, domain=0.6:4.6, samples=150]
    {0.42*(x-0.6)^2 - 0.06*(x-0.6)^3 + 0.5};
  % corde
  \addplot[bmtGris, line width=0.9pt, domain=0.6:4.6, samples=2]
    {0.5 + 0.72*(x-0.6)};
  % tangente au point c
  \addplot[bmtLaterite, dashed, line width=0.95pt, domain=0.75:3.25, samples=2]
    {0.9508 + 0.72*(x-1.7315)};
  \filldraw[bmtLaterite] (axis cs:1.7315,0.9508) circle (1.9pt);
  \filldraw[bmtGris] (axis cs:0.6,0.5) circle (1.7pt);
  \filldraw[bmtGris] (axis cs:4.6,3.38) circle (1.7pt);
  \node[bmtGris, font=\footnotesize, anchor=south east] at (axis cs:4.45,3.44) {la corde};
  \node[bmtLaterite, font=\footnotesize, anchor=north west] at (axis cs:2.6,1.55) {la tangente en $c$};
\end{axis}
\end{tikzpicture}
\end{center}

Joignez les deux extrémités de l'arc par une corde : sa pente vaut
$\dfrac{f(b)-f(a)}{b-a}$. Le théorème affirme qu'entre $a$ et $b$ il existe au
moins un point où la tangente à la courbe est \emph{parallèle} à cette corde.
Faites glisser mentalement la corde vers le haut sans changer sa direction :
elle finit par toucher la courbe, et ce point de contact est le $c$ cherché.
\end{visualisation}

\begin{theoreme}[inégalité des accroissements finis]
Soit $f$ continue sur $\intervalle{a}{b}$, dérivable sur $\intervalleo{a}{b}$.
S'il existe deux réels $m$ et $M$ tels que $m \leq f'(x) \leq M$ pour tout
$x \in \intervalleo{a}{b}$, alors
\[
  m(b-a) \ \leq\ f(b)-f(a) \ \leq\ M(b-a) .
\]
En particulier, si $\abs{f'(x)} \leq k$ sur $\intervalleo{a}{b}$, alors
\[
  \abs{f(b)-f(a)} \ \leq\ k\,\abs{b-a} .
\]
\end{theoreme}

\begin{demonstration}
Le théorème des accroissements finis fournit un réel
$c \in \intervalleo{a}{b}$ tel que $f(b)-f(a) = (b-a)f'(c)$. Comme
$b-a > 0$, l'encadrement $m \leq f'(c) \leq M$ donne, après multiplication par
$b-a$,
$m(b-a) \leq (b-a)f'(c) \leq M(b-a)$, c'est-à-dire l'encadrement annoncé.
Pour la seconde forme, $\abs{f'(c)} \leq k$ entraîne
$\abs{f(b)-f(a)} = \abs{b-a}\,\abs{f'(c)} \leq k\abs{b-a}$.
\end{demonstration}

\begin{methode}{majorer un écart sans calculer la fonction}
L'inégalité des accroissements finis sert chaque fois qu'on doit majorer
$\abs{f(x)-f(y)}$ sans pouvoir — ou sans vouloir — calculer ces deux valeurs.
La marche à suivre est toujours la même.

\begin{enumerate}[label=\textbf{\arabic*.}, leftmargin=*, itemsep=0.3em]
  \item Identifier l'intervalle sur lequel on travaille et vérifier que $f$ y
        est continue et dérivable.
  \item Calculer $f'$ et \emph{majorer $\abs{f'}$} sur cet intervalle : c'est
        la seule étape qui demande du soin, et c'est celle qui rapporte les
        points.
  \item Appliquer l'inégalité avec la constante $k$ obtenue.
\end{enumerate}

C'est exactement le mécanisme qui, appliqué à une suite définie par
$u_{n+1} = f(u_{n})$, donnera $\abs{u_{n+1}-\alpha} \leq k\abs{u_{n}-\alpha}$
puis, par récurrence, $\abs{u_{n}-\alpha} \leq k^{n}\abs{u_{0}-\alpha}$.
\end{methode}

\begin{exemple}[un encadrement obtenu sans calculatrice]
Montrons que $\abs{\sin x - \sin y} \leq \abs{x-y}$ pour tous réels $x$ et $y$.

La fonction sinus est continue et dérivable sur $\R$, de dérivée
$x \mapsto \cos x$. Or $\abs{\cos t} \leq 1$ pour tout réel $t$.

L'inégalité des accroissements finis, appliquée sur l'intervalle d'extrémités
$x$ et $y$ avec $k = 1$, donne immédiatement
$\abs{\sin x - \sin y} \leq \abs{x-y}$.

\smallskip
En particulier, avec $y = 0$ : $\abs{\sin x} \leq \abs{x}$ pour tout réel $x$
— une inégalité qu'on utilisera plus d'une fois.
\end{exemple}

\begin{entrainement}[accroissements finis]
\exo[\niveau{1}] Soit $f(x) = x^{2}$ sur $\intervalle{1}{3}$. Déterminer le
réel $c$ dont le théorème des accroissements finis garantit l'existence.

\exo[\niveau{2}] Montrer que pour tous réels $x$ et $y$ positifs,
$\abs{\sqrt{x}-\sqrt{y}} \leq \dfrac{\abs{x-y}}{2}$ dès que $x \geq 1$ et
$y \geq 1$.

\exo[\niveau{2}] Soit $f$ dérivable sur $\intervalle{0}{4}$ avec $f(0) = 2$ et
$\abs{f'(x)} \leq 3$. Entre quelles bornes est nécessairement compris $f(4)$ ?

\exo[\niveau{3}] Soit $f(x) = \dfrac{x^{2}+2}{4}$. Montrer que pour tout
$x \in \intervalle{0}{1}$, $f(x) \in \intervalle{0}{1}$ et
$\abs{f'(x)} \leq \dfrac{1}{2}$.
\end{entrainement}

\begin{corrige}
\textbf{1.} $f(3)-f(1) = 8 = 2 \times f'(c) = 2 \times 2c$, d'où $c = 2$, qui
appartient bien à $\intervalleo{1}{3}$.

\textbf{2.} Sur $\left[1 ; +\infty\right[$,
$\left(\sqrt{t}\right)' = \dfrac{1}{2\sqrt{t}} \leq \dfrac{1}{2}$ puisque
$\sqrt{t} \geq 1$. L'inégalité des accroissements finis avec $k = \tfrac{1}{2}$
donne le résultat.

\textbf{3.} $\abs{f(4)-f(0)} \leq 3 \times 4 = 12$, donc
$-10 \leq f(4) \leq 14$.

\textbf{4.} Pour $x \in \intervalle{0}{1}$, $x^{2} \in \intervalle{0}{1}$ donc
$f(x) \in \intervalle{\tfrac{1}{2}}{\tfrac{3}{4}} \subset \intervalle{0}{1}$.
Et $f'(x) = \dfrac{x}{2}$, donc $\abs{f'(x)} \leq \dfrac{1}{2}$ sur cet
intervalle.
\end{corrige}

% ===========================================================================
% ===========================================================================
\begin{annales}
Les énoncés qui suivent sont repris de sujets réellement posés ; leur
provenance est indiquée à chaque fois. Tous portent sur des fonctions
polynômes, rationnelles ou irrationnelles : rien n'y suppose une notion non
encore étudiée.

\exoann{Baccalauréat probatoire, Cameroun, série D, session 2022 — exercice 4}
Soit $f$ définie sur $\intervallefo{0}{+\infty}$ par
$f(x) = \dfrac{3x}{3+4x}$. Calculer $f'(x)$, en déduire le sens de variation
de $f$, puis dresser son tableau de variation.

\exoann{Baccalauréat probatoire, Cameroun, séries D et TI, session 2023 — exercice 4}
Soit $f$ définie par $f(x) = x+\dfrac{1}{x+1}$. Calculer $f'(x)$, étudier son
signe et en déduire les variations de $f$ sur chacun des intervalles où elle
est définie.

\exoann{Baccalauréat probatoire, Cameroun, série IH, session 2023 — problème}
Soit $f$ la fonction définie sur $\intervalle{-4}{1}$ par
$f(x) = ax^{2}+bx-3$, dont la courbe passe par $A(-2\,;-5)$ et
$B(-1\,;-6)$. Déterminer les réels $a$ et $b$.

\exoann{Baccalauréat probatoire, Cameroun, série IH, session 2023 — problème, suite}
Avec les valeurs de $a$ et $b$ trouvées, déterminer une équation de la
tangente à la courbe de $f$ au point $A(-2\,;-5)$.

\exoann{D'après le baccalauréat probatoire, Cameroun, série C, session 2021 — exercice 3}
Soit $f$ définie sur $\R\setminus\{1\}$ par $f(x) = \dfrac{x^{2}-4}{1-x}$.
Vérifier que $f'(x) = \dfrac{-x^{2}+2x-4}{(1-x)^{2}}$, puis en déduire le sens
de variation de $f$ sur chacun des intervalles où elle est définie.

\exoann{D'après le baccalauréat probatoire, Cameroun, série C, session 2022 — exercice 2}
Soit $f$ définie sur $\R$ par
$f(x) = \dfrac{1}{3}x^{3}+\dfrac{3}{2}x^{2}-4x-\dfrac{31}{6}$. Calculer
$f'(x)$, étudier son signe, puis dresser le tableau de variation de $f$ en
précisant les deux extremums locaux.

\exoann{D'après le baccalauréat probatoire, Cameroun, série C, session 2023 — exercice 4}
Soit $f$ définie par $f(x) = \dfrac{2x^{2}-6x+3}{2x-3}$. Calculer $f'(x)$ et
en déduire le sens de variation de $f$ sur chacun des intervalles où elle est
définie.

\exoann{D'après le baccalauréat probatoire, Cameroun, série C, session 2024 — exercice 3}
Soit $f$ définie par $f(x) = \dfrac{4x^{2}-4x-2}{2x-3}$.
\begin{enumerate}[label=\textbf{\alph*)}, leftmargin=*, itemsep=0.15em]
  \item Vérifier que $f(x) = 2x+1+\dfrac{1}{2x-3}$.
  \item Calculer $f'(x)$, puis résoudre l'inéquation $f'(x) \leq 0$.
\end{enumerate}

\exoann{D'après le baccalauréat probatoire, Cameroun, série C, session 2024 — exercice 3, suite}
Avec la même fonction $f$, démontrer que le point
$\Omega\left(\frac32\,;4\right)$ est centre de symétrie de la courbe de $f$.

\medskip
\bmtsource{Les autres énoncés d'annales portant sur la dérivabilité font
intervenir le logarithme ou l'exponentielle. Ils figurent donc aux
chapitres 5 et 6, où ils sont signalés par la mention « dérivabilité ».}
\end{annales}

\begin{corrige}[annales]
\textbf{1.} $f'(x) = \dfrac{3(3+4x)-3x \times 4}{(3+4x)^{2}}
= \dfrac{9}{(3+4x)^{2}}$, strictement positive sur
$\intervallefo{0}{+\infty}$ : $f$ y est strictement croissante, de
$f(0) = 0$ vers $\frac34$, valeur qu'elle n'atteint jamais.

\textbf{2.} $f'(x) = 1-\dfrac{1}{(x+1)^{2}}
= \dfrac{(x+1)^{2}-1}{(x+1)^{2}} = \dfrac{x(x+2)}{(x+1)^{2}}$. Le
dénominateur est positif, donc $f'$ a le signe de $x(x+2)$ : positive sur
$\left]-\infty\,;-2\right]$, négative sur $\intervalle{-2}{0}$ privé de
$-1$, positive sur $\intervallefo{0}{+\infty}$. La fonction croît, atteint un
maximum local $f(-2) = -3$, décroît jusqu'à l'asymptote verticale, décroît
encore, atteint un minimum local $f(0) = 1$, puis croît.

\textbf{3.} Le point $A$ donne $4a-2b-3 = -5$, soit $2a-b = -1$ ; le point
$B$ donne $a-b-3 = -6$, soit $a-b = -3$. En soustrayant, $a = 2$, puis
$b = 5$. Ainsi $f(x) = 2x^{2}+5x-3$, ce que confirme le calcul de $f(-2)$ et
de $f(-1)$.

\textbf{4.} $f'(x) = 4x+5$, donc $f'(-2) = -3$. La tangente en $A$ a pour
équation $y = -3\left(x+2\right)-5$, c'est-à-dire $y = -3x-11$.

\textbf{5.} En dérivant le quotient,
\[
  f'(x) = \frac{2x(1-x)-\left(x^{2}-4\right)\times(-1)}{(1-x)^{2}}
  = \frac{2x-2x^{2}+x^{2}-4}{(1-x)^{2}}
  = \frac{-x^{2}+2x-4}{(1-x)^{2}} .
\]
Le trinôme $-x^{2}+2x-4$ a pour discriminant $4-16 = -12 < 0$ : il garde le
signe de son coefficient dominant, donc il est strictement négatif. Le
dénominateur étant strictement positif, $f'$ est strictement négative :
$f$ décroît strictement sur $\left]-\infty\,;1\right[$ et sur
$\intervalleo{1}{+\infty}$.

\textbf{6.} $f'(x) = x^{2}+3x-4 = (x+4)(x-1)$, positive sur
$\left]-\infty\,;-4\right]$, négative sur $\intervalle{-4}{1}$, positive sur
$\intervallefo{1}{+\infty}$. La fonction croît donc jusqu'au maximum local
$f(-4) = \frac{27}{2}$, décroît jusqu'au minimum local
$f(1) = -\frac{22}{3}$, puis croît de nouveau ; les limites aux infinis sont
celles de $\frac{x^{3}}{3}$, soit $-\infty$ et $+\infty$.

\textbf{7.} En dérivant le quotient,
\[
\begin{aligned}
  f'(x) &= \frac{(4x-6)(2x-3)-2\left(2x^{2}-6x+3\right)}{(2x-3)^{2}}\\
  &= \frac{8x^{2}-24x+18-4x^{2}+12x-6}{(2x-3)^{2}}\\
  &= \frac{4\left(x^{2}-3x+3\right)}{(2x-3)^{2}} .
\end{aligned}
\]
Le trinôme $x^{2}-3x+3$ a pour discriminant $9-12 = -3<0$ : il est toujours
strictement positif. Donc $f'>0$ et $f$ croît strictement sur
$\left]-\infty\,;\frac32\right[$ comme sur
$\left]\frac32\,;+\infty\right[$ — sans que ces deux croissances se
raccordent, l'asymptote verticale les séparant.

\textbf{8. a)} On développe le membre de droite :
$(2x+1)(2x-3)+1 = 4x^{2}-4x-3+1 = 4x^{2}-4x-2$, qui est bien le numérateur de
$f$. D'où $f(x) = 2x+1+\dfrac{1}{2x-3}$.
\textbf{b)} $f'(x) = 2-\dfrac{2}{(2x-3)^{2}}$, et
\[
  f'(x) \leq 0
  \iff (2x-3)^{2} \leq 1
  \iff -1 \leq 2x-3 \leq 1
  \iff 1 \leq x \leq 2 .
\]
La valeur $\frac32$ étant exclue, $f'$ est négative sur
$\intervallefo{1}{\frac32}$ et sur $\left]\frac32\,;2\right]$, positive
ailleurs.

\textbf{9.} On teste la symétrie par le calcul de $f(\Omega_{x}+h)+f(\Omega_{x}-h)$.
Pour $h \neq 0$ tel que $\frac32 \pm h$ soit dans l'ensemble de définition,
\[\begin{aligned}
&f\!\left(\tfrac32+h\right)+f\!\left(\tfrac32-h\right)\\
  &\quad = \left(4+2h+\frac{1}{2h}\right)+\left(4-2h-\frac{1}{2h}\right)
  \\ &\quad = 8 = 2\times4 .
\end{aligned}\]
La moyenne des deux ordonnées vaut donc toujours $4$, ordonnée de $\Omega$ :
le point $\Omega\left(\frac32\,;4\right)$ est centre de symétrie de la courbe.
\end{corrige}

\begin{jecherche}[la descente d'un col]
Un cycliste descend un col. On note $d(t)$ la distance parcourue, en
kilomètres, $t$ heures après le départ. La fonction $d$ est dérivable sur
$\intervalle{0}{2}$, et $d'(t)$ représente la vitesse instantanée, en
kilomètres par heure.

On sait que $d(0) = 0$ et $d(2) = 84$.

\begin{enumerate}[label=\textbf{\arabic*.}, leftmargin=*, itemsep=0.3em]
  \item Justifier qu'il existe un instant où la vitesse instantanée du
        cycliste est exactement égale à sa vitesse moyenne.
  \item Le règlement de l'épreuve interdit de dépasser $60$ km/h. Le cycliste
        affirme l'avoir respecté. Que pensez-vous de cette affirmation ?
  \item On suppose maintenant que la vitesse est restée comprise entre
        $30$ et $50$ km/h sur toute la descente. Entre quelles valeurs la
        distance parcourue en deux heures peut-elle se situer ? Cette
        hypothèse est-elle compatible avec les données ?
  \item Reprendre la question 2 en supposant seulement que la vitesse a été
        mesurée à $40$ km/h au départ et à $45$ km/h à l'arrivée. Peut-on
        encore conclure ?
\end{enumerate}
\end{jecherche}

\begin{corrige}[la descente d'un col]
\textbf{1.} $d$ est continue sur $\intervalle{0}{2}$ et dérivable sur
$\intervalleo{0}{2}$. Le théorème des accroissements finis fournit
$c \in \intervalleo{0}{2}$ tel que $d(2)-d(0) = 2\,d'(c)$, soit
$d'(c) = \dfrac{84}{2} = 42$. Il existe donc un instant où la vitesse
instantanée vaut exactement $42$ km/h, la vitesse moyenne.

\textbf{2.} La vitesse moyenne vaut $42$ km/h, bien en dessous de la limite.
On ne peut donc rien reprocher au cycliste — mais on ne peut rien confirmer non
plus : une moyenne ne majore pas les valeurs instantanées. Il a pu rouler à
$70$ km/h pendant dix minutes et à $30$ le reste du temps. La moyenne est
compatible avec le respect du règlement comme avec son infraction.

\textbf{3.} Si $30 \leq d'(t) \leq 50$ sur $\intervalleo{0}{2}$, l'inégalité
des accroissements finis donne
$30 \times 2 \leq d(2)-d(0) \leq 50 \times 2$, soit
$60 \leq d(2) \leq 100$. La valeur $84$ appartient à cet intervalle :
l'hypothèse est compatible avec les données.

\textbf{4.} Non. Connaître $d'$ en deux points ne dit rien de $d'$ ailleurs :
l'inégalité des accroissements finis exige un encadrement sur \emph{tout}
l'intervalle. Une fonction peut valoir $40$ en $0$, $45$ en $2$, et culminer à
$90$ entre les deux.
\end{corrige}


% ===========================================================================
\begin{jemeteste}
Répondre par vrai ou faux, en justifiant en une ligne.

\begin{enumerate}[label=\textbf{\arabic*.}, leftmargin=*, itemsep=0.28em]
  \item Toute fonction continue sur $\R$ est dérivable sur $\R$.
  \item Si $f'_{g}(a)$ et $f'_{d}(a)$ existent, alors $f$ est dérivable en $a$.
  \item Si $f''(a) = 0$, alors la courbe de $f$ admet un point d'inflexion
        d'abscisse $a$.
  \item Si $f$ est convexe sur $I$, sa courbe est au-dessus de chacune de ses
        tangentes sur $I$.
  \item Si $\abs{f'} \leq 2$ sur $\intervalle{0}{5}$, alors
        $\abs{f(5)-f(0)} \leq 10$.
  \item Si $f$ est dérivable et strictement croissante sur $\R$, alors
        $f'(x) > 0$ pour tout $x$.
\end{enumerate}
\end{jemeteste}

\begin{corrige}[je me teste]
\textbf{1.} Faux — $x \mapsto \abs{x}$ n'est pas dérivable en $0$.
\textbf{2.} Faux — il faut de plus qu'ils soient égaux.
\textbf{3.} Faux — il faut un changement de signe ; contre-exemple
$f(x) = x^{4}$.
\textbf{4.} Vrai.
\textbf{5.} Vrai — inégalité des accroissements finis avec $k = 2$.
\textbf{6.} Faux — $f(x) = x^{3}$ est strictement croissante et $f'(0) = 0$.
\end{corrige}

% ===========================================================================
\begin{commeaubac}
\bmtsource{Exercice au format de l'épreuve}
Soit $f$ la fonction définie sur $\R$ par $f(x) = \dfrac{x^{3}}{x^{2}+1}$.

\begin{enumerate}[label=\textbf{\arabic*.}, leftmargin=*, itemsep=0.25em]
  \item Étudier la parité de $f$ et en déduire un élément de symétrie de sa
        courbe. \hfill\textup{(0,5 pt)}
  \item Calculer $f'(x)$ et montrer que $f$ est strictement croissante
        sur $\R$. \hfill\textup{(1 pt)}
  \item Montrer que la droite d'équation $y = x$ est asymptote à la courbe en
        $+\infty$ et en $-\infty$, et étudier la position relative de la
        courbe et de cette droite. \hfill\textup{(1 pt)}
  \item Calculer $f''(x)$ et déterminer les points d'inflexion de la courbe.
        \hfill\textup{(1,5 pt)}
  \item Montrer que pour tous réels $x$ et $y$,
        $\abs{f(x)-f(y)} \leq \dfrac{9}{8}\abs{x-y}$.
        \hfill\textup{(1 pt)}
\end{enumerate}
\end{commeaubac}

\begin{corrige}[exercice au format de l'épreuve]
\textbf{1.} $f(-x) = \dfrac{-x^{3}}{x^{2}+1} = -f(x)$ : $f$ est impaire, sa
courbe est symétrique par rapport à l'origine.

\textbf{2.} $f'(x) = \dfrac{3x^{2}(x^{2}+1) - x^{3}\cdot 2x}{(x^{2}+1)^{2}}
= \dfrac{x^{4}+3x^{2}}{(x^{2}+1)^{2}} = \dfrac{x^{2}(x^{2}+3)}{(x^{2}+1)^{2}} \geq 0$,
nulle seulement en $0$ : $f$ est strictement croissante sur $\R$.

\textbf{3.} $f(x) - x = \dfrac{x^{3}-x(x^{2}+1)}{x^{2}+1} = \dfrac{-x}{x^{2}+1} \to 0$
aux deux infinis. Le signe de $f(x)-x$ est celui de $-x$ : la courbe est
au-dessus de la droite pour $x<0$, au-dessous pour $x>0$.

\textbf{4.} En dérivant $f'$ on obtient
$f''(x) = \dfrac{2x(3-x^{2})}{(x^{2}+1)^{3}}$, qui s'annule en changeant de
signe pour $x = -\sqrt{3}$, $x = 0$ et $x = \sqrt{3}$ : trois points
d'inflexion.

\textbf{5.} On étudie $f'$ : en posant $t = x^{2} \geq 0$,
$f'(x) = \dfrac{t(t+3)}{(t+1)^{2}}$. Cette expression est maximale pour $t=3$,
où elle vaut $\dfrac{3 \times 6}{16} = \dfrac{9}{8}$. Donc
$0 \leq f'(x) \leq \dfrac{9}{8}$ sur $\R$, et l'inégalité des accroissements
finis donne le résultat.
\end{corrige}

% ===========================================================================
\begin{concours}
\bmtsource{D'après le concours d'entrée de l'ENI, Fianarantsoa}
Soit $f$ la fonction définie sur $\R$ par $f(x) = \dfrac{x^{2}+3x+4m}{x^{2}+(5m+1)x+3}$,
où $m$ est un paramètre réel.

\begin{enumerate}[label=\textbf{\arabic*.}, leftmargin=*, itemsep=0.25em]
  \item Pour $m = 0$, montrer que $f$ est dérivable sur $\R$ et calculer
        $f'(x)$.
  \item Toujours pour $m = 0$, déterminer le nombre de points de la courbe où
        la tangente est horizontale.
  \item Pour quelles valeurs de $m$ la tangente à la courbe au point
        d'abscisse $0$ est-elle parallèle à la droite d'équation $y = x$ ?
\end{enumerate}
\end{concours}

\begin{corrige}[vers le concours]
\textbf{1.} Pour $m = 0$, le dénominateur vaut $x^{2}+x+3$, de discriminant
$1-12 = -11 < 0$ : il ne s'annule jamais, donc $f_{0}$ est dérivable sur $\R$
comme quotient de fonctions polynômes. En dérivant,
\[
  f_{0}'(x) = \frac{(2x+3)(x^{2}+x+3)-(x^{2}+3x)(2x+1)}{(x^{2}+x+3)^{2}}
            = \frac{-2x^{2}+6x+9}{(x^{2}+x+3)^{2}} .
\]

\textbf{2.} Une tangente est horizontale lorsque $f_{0}'(x) = 0$, donc lorsque
$-2x^{2}+6x+9 = 0$. Le discriminant vaut $36+72 = 108 > 0$ : deux racines
réelles, donc \emph{deux} points à tangente horizontale.

\textbf{3.} La tangente en $0$ est parallèle à $y = x$ lorsque $f_{m}'(0) = 1$.
Or
\[
  f_{m}'(0) = \frac{3 \times 3 - 4m(5m+1)}{3^{2}} = \frac{9-20m^{2}-4m}{9},
\]
et l'égalité $f_{m}'(0) = 1$ donne $-4m(5m+1) = 0$, soit $m = 0$ ou
$m = -\dfrac{1}{5}$. Ces deux valeurs conviennent : pour $m = -\frac15$ le
dénominateur devient $x^{2}+3$, qui ne s'annule pas non plus.
\end{corrige}


% ===========================================================================
\begin{devoir}[2 heures]
Trois exercices indépendants, notés sur $20$. Le devoir porte sur ce chapitre
et sur celui qui le précède. Les énoncés sont repris de sessions réelles et
assemblés ici en épreuve ; leur provenance est indiquée à chaque fois.

\exods{1}{D'après le baccalauréat probatoire, Cameroun, série C, session 2020 — exercice 3, partie I}{5 points}
Le plan est rapporté à un repère orthonormé. On considère les points
$A(0\,;6)$ et $B(-2\,;-4)$, et l'ensemble $(\Gamma)$ des points $M$ du plan
vérifiant $AM^{2}+BM^{2} = 102$.
\begin{enumerate}[label=\textbf{\arabic*.}, leftmargin=*, itemsep=0.2em]
  \item Déterminer une équation cartésienne de $(\Gamma)$.
        \hfill\textup{(2 pts)}
  \item En déduire la nature de $(\Gamma)$ et ses éléments caractéristiques.
        \hfill\textup{(1,5 pt)}
  \item Vérifier que le centre de $(\Gamma)$ est le milieu $I$ de $[AB]$, et
        retrouver le rayon à l'aide de l'égalité
        $AM^{2}+BM^{2} = 2\,IM^{2}+\dfrac{AB^{2}}{2}$.
        \hfill\textup{(1,5 pt)}
\end{enumerate}

\exods{2}{D'après le baccalauréat probatoire, Cameroun, série C, session 2021 — exercice 3}{7 points}
Soit $f$ la fonction définie sur $\R\setminus\{1\}$ par
$f(x) = \dfrac{x^{2}-4}{1-x}$, de courbe $(C)$.
\begin{enumerate}[label=\textbf{\arabic*.}, leftmargin=*, itemsep=0.2em]
  \item Vérifier que $f'(x) = \dfrac{-x^{2}+2x-4}{(1-x)^{2}}$.
        \hfill\textup{(2 pts)}
  \item Étudier le signe de $f'(x)$, puis donner le sens de variation de $f$
        sur chacun des intervalles où elle est définie.
        \hfill\textup{(2 pts)}
  \item Déterminer une équation de la tangente $(T_{0})$ à $(C)$ au point
        d'abscisse $0$. \hfill\textup{(1,5 pt)}
  \item Étudier la position de $(C)$ par rapport à $(T_{0})$ sur
        $\left]-\infty\,;1\right[$. \hfill\textup{(1,5 pt)}
\end{enumerate}

\exods{3}{D'après le baccalauréat probatoire, Cameroun, série D, session 2024 — exercice II}{8 points}
Soit $f$ la fonction définie sur $\R$ par $f(x) = -x^{3}+3x^{2}$.
\begin{enumerate}[label=\textbf{\arabic*.}, leftmargin=*, itemsep=0.2em]
  \item Calculer $f'(x)$ et étudier son signe. \hfill\textup{(1,5 pt)}
  \item Dresser le tableau de variation de $f$ en précisant les deux
        extremums locaux. \hfill\textup{(2 pts)}
  \item Calculer $f''(x)$ et déterminer le point d'inflexion de la courbe de
        $f$. \hfill\textup{(2 pts)}
  \item Déterminer une équation de la tangente à la courbe en ce point
        d'inflexion. \hfill\textup{(1,5 pt)}
  \item On pose $g(x) = f(x-1)$. En déduire, sans nouveau calcul de dérivée,
        le tableau de variation de $g$. \hfill\textup{(1 pt)}
\end{enumerate}
\end{devoir}

\begin{corrige}[devoir surveillé]
\textbf{Exercice 1. 1.} Pour $M(x\,;y)$,
$AM^{2} = x^{2}+(y-6)^{2}$ et $BM^{2} = (x+2)^{2}+(y+4)^{2}$. La somme vaut
$2x^{2}+2y^{2}+4x-4y+56$, et l'égalité $AM^{2}+BM^{2} = 102$ devient, après
division par $2$,
\[
  x^{2}+y^{2}+2x-2y-23 = 0,
  \qquad\text{soit}\qquad
  (x+1)^{2}+(y-1)^{2} = 25 .
\]

\textbf{2.} C'est l'équation d'un cercle de centre $I(-1\,;1)$ et de rayon
$5$.

\textbf{3.} Le milieu de $[AB]$ a pour coordonnées
$\left(\frac{0-2}{2}\,;\frac{6-4}{2}\right) = (-1\,;1)$ : c'est bien $I$. Par
ailleurs $AB^{2} = (-2)^{2}+(-10)^{2} = 104$, donc l'égalité de la médiane
donne $102 = 2\,IM^{2}+52$, d'où $IM^{2} = 25$ et $IM = 5$. On retrouve le
cercle de centre $I$ et de rayon $5$ sans passer par les coordonnées.

\medskip
\textbf{Exercice 2. 1.} En dérivant le quotient,
\[
  f'(x) = \frac{2x(1-x)-\left(x^{2}-4\right)\times(-1)}{(1-x)^{2}}
  = \frac{2x-2x^{2}+x^{2}-4}{(1-x)^{2}}
  = \frac{-x^{2}+2x-4}{(1-x)^{2}} .
\]

\textbf{2.} Le trinôme $-x^{2}+2x-4$ a pour discriminant $4-16 = -12<0$ : il
ne s'annule jamais et garde le signe de son coefficient dominant, donc il est
strictement négatif. Le dénominateur étant strictement positif, $f'<0$ :
$f$ décroît strictement sur $\left]-\infty\,;1\right[$ et sur
$\intervalleo{1}{+\infty}$.

\textbf{3.} $f(0) = \dfrac{-4}{1} = -4$ et
$f'(0) = \dfrac{-4}{1} = -4$. La tangente a donc pour équation
$y = -4(x-0)-4$, c'est-à-dire $y = -4x-4$.

\textbf{4.} On étudie la différence :
\[
  f(x)-(-4x-4) = \frac{x^{2}-4}{1-x}+4x+4
  = \frac{x^{2}-4+(4x+4)(1-x)}{1-x}
  = \frac{-3x^{2}}{1-x} .
\]
Sur $\left]-\infty\,;1\right[$, le dénominateur est strictement positif et le
numérateur négatif ou nul : la différence est négative, nulle seulement en
$0$. La courbe est donc au-dessous de sa tangente sur tout cet intervalle,
qu'elle touche au seul point de contact.

\medskip
\textbf{Exercice 3. 1.} $f'(x) = -3x^{2}+6x = -3x(x-2)$, négatif à
l'extérieur de $\intervalle{0}{2}$ et positif à l'intérieur.

\textbf{2.} La fonction décroît sur $\left]-\infty\,;0\right]$, croît sur
$\intervalle{0}{2}$, décroît sur $\intervallefo{2}{+\infty}$. Elle admet un
minimum local $f(0) = 0$ et un maximum local $f(2) = -8+12 = 4$. Les limites
sont $+\infty$ en $-\infty$ et $-\infty$ en $+\infty$.

\textbf{3.} $f''(x) = -6x+6$, qui s'annule en $x = 1$ en changeant de signe :
la courbe est convexe avant $1$, concave après. Le point d'inflexion est
$\left(1\,;f(1)\right) = (1\,;2)$.

\textbf{4.} $f'(1) = -3+6 = 3$, donc la tangente en ce point a pour équation
$y = 3(x-1)+2$, soit $y = 3x-1$.

\textbf{5.} La courbe de $g$ se déduit de celle de $f$ par la translation qui
déplace chaque point d'une unité vers la droite : les variations sont les
mêmes, décalées de $1$. Ainsi $g$ décroît sur $\left]-\infty\,;1\right]$,
croît sur $\intervalle{1}{3}$, décroît sur $\intervallefo{3}{+\infty}$, avec
un minimum local $g(1) = 0$ et un maximum local $g(3) = 4$.
\end{corrige}
